% from ega0-17-fr.tex
\oldpage[0\textsubscript{IV-1}]{50}

\begin{corollary}[17.3.5]
\label{0.17.3.5-fr}
\begin{enumerate}
  \item[\rm(i)] Soit $A$ un anneau local régulier de dimension $n$~; pour
  qu'un $A$-module $M\neq0$ de type fini soit libre, il faut et il suffit que
  $M$ soit un $A$-module de Cohen-Macaulay de dimension $n$.
  \item[\rm(ii)] Soient $A$ un anneau local régulier, $B$ un anneau local,
  $\rho:A\to B$ un homomorphisme local faisant de $B$ un $A$-module de type
  fini. Pour que $B$ soit un $A$-module libre, il faut et il suffit que $B$
  soit un anneau de Cohen-Macaulay et que $\rho$ soit injectif (ou, ce qui
  revient au même (16.1.5), que $\dim(B)=\dim(A)$).
\end{enumerate}
\end{corollary}

\begin{proof}
(i) Cela résulte de (17.3.4.1) et du fait que pour un $A$-module de type fini,
il revient au même de dire que ce module est projectif ou libre (Bourbaki,
\emph{Alg. comm.}, chap.~II, \S~3, n\textsuperscript{o}~2, cor.~2 de la
prop.~5)~; les $A$-modules libres de type fini $M$ sont donc caractérisés par
la relation $\operatorname{dim.proj}(M)=0$ (17.2.2, (i)).

(ii) Dire que $B$ est un anneau de Cohen-Macaulay équivaut à dire que $B$ est
un $A$-module de Cohen-Macaulay (16.5.3), donc il suffit d'appliquer (i),
puisque $\dim(B)=\dim(A)$ (16.1.5).
\end{proof}

\begin{env}[17.3.6]
\label{0.17.3.6-fr}
On dit qu'un anneau noethérien $A$ est \emph{régulier} si pour tout idéal
premier $\mathfrak{p}$ de $A$, l'anneau local $A_{\mathfrak{p}}$ est
régulier~; lorsque $A$ lui-même est un anneau local, il résulte de (17.3.2)
que cette définition est équivalente à celle de (17.1.2). Pour que $A$ soit
régulier, il suffit, en vertu de (17.3.2), que $A_{\mathfrak{m}}$ soit
régulier pour tout idéal maximal $\mathfrak{m}$ de $A$. En outre, il résulte
aussitôt de cette définition que pour toute partie multiplicative $S$ de
$A$, $S^{-1}A$ est régulier.
\end{env}

\begin{proposition}[17.3.7]
\label{0.17.3.7-fr}
Si $A$ est un anneau noethérien régulier, tout anneau de polynômes
$A[T_1,\ldots,T_n]$ est régulier.
\end{proposition}

\begin{proof}
Il suffit évidemment de prouver que l'anneau de polynômes $B=A[T]$ est
régulier~; comme $B$ est un $A$-module libre, pour tout idéal premier
$\mathfrak{q}$ de $B$, $B_{\mathfrak{q}}$ est un
$A_{\mathfrak{p}}$-module plat, où $\mathfrak{p}=\mathfrak{q}\cap A$
($0_{\mathrm I}$, 6.3.1)~; il suffit donc, en vertu de (17.1.10), de prouver
que $B_{\mathfrak{q}}/\mathfrak{p}B_{\mathfrak{q}}$ est régulier, et comme
cet anneau est un anneau local en un idéal premier de
$A_{\mathfrak{p}}[T]/\mathfrak{p}A_{\mathfrak{p}}[T]=k[T]$ (où $k$ est le
corps résiduel $A_{\mathfrak{p}}/\mathfrak{p}A_{\mathfrak{p}}$), il suffit
de prouver que $C=k[T]$ est régulier~; or, $C$ étant principal, les anneaux
locaux en les idéaux premiers de $C$ sont des anneaux de valuation discrète
ou un corps (pour l'idéal $(0)$), donc réguliers (17.1.4), ce qui achève la
démonstration.
\end{proof}

\begin{corollary}[17.3.8]
\label{0.17.3.8-fr}
Si $A$ est un anneau régulier, tout anneau de séries formelles
$A[[T_1,\ldots,T_n]]$ est régulier.
\end{corollary}

\begin{proof}
Soit $\mathfrak{J}$ l'idéal engendré par les $T_i$ dans l'anneau de polynômes
$A[T_1,\ldots,T_n]$~; comme ce dernier est régulier par (17.3.7), et que
$A[[T_1,\ldots,T_n]]$ est le complété de $A[T_1,\ldots,T_n]$ pour la
topologie $\mathfrak{J}$-préadique (Bourbaki, \emph{Alg. comm.}, chap.~III,
\S~2, n\textsuperscript{o}~12, \emph{Exemple 1}), la conclusion résulte du
\end{proof}

\begin{lemma}[17.3.8.1]
\label{0.17.3.8.1-fr}
Soient $A$ un anneau régulier, $\mathfrak{J}$ un idéal de $A$,
$\widehat{A}$ le séparé complété de $A$ pour la topologie
$\mathfrak{J}$-préadique. Alors $\widehat{A}$ est régulier.
\end{lemma}

\begin{proof}
Il suffit en effet (17.3.6) de voir que pour tout idéal maximal
$\mathfrak{n}$ de $\widehat{A}$, l'anneau local
$\widehat{A}_{\mathfrak{n}}$ est régulier~; on sait (Bourbaki,
\emph{Alg. comm.}, chap.~III, \S~3, n\textsuperscript{o}~4, prop.~8) que
$\mathfrak{n}$
